\chapter{Strings, Gravity, Compactification, and T-Duality}

\lectureoverview{A point particle can move consistently on essentially any
smooth geometry, but a quantum string cannot.  Its infinitely many
zero-point oscillations probe the target space at successively shorter
scales, and the effective geometry changes as those oscillations are
restored.  Requiring that this process have a fixed point gives
\(R_{\mu\nu}=0\), the vacuum Einstein equation.  We then turn from this
consistency condition to compactification.  A circle supplies both
Kaluza--Klein momentum states and string winding states; exchanging them
identifies a circle of radius \(R\) with one of radius
\(\alpha'/R\).}

\section{Consistency as a Selection Principle}

There are two questions to answer.  Why does gravity occur in string theory,
and why is string theory far less permissive than point-particle mechanics?
They are really one question.  A background geometry is acceptable only if
a string can propagate on it without retaining an arbitrary dependence on
how finely we resolve the string.

Begin with the control problem: a point particle constrained to a smooth
curved surface.  Classically it follows a geodesic, the intrinsic analogue
of a straight line.  One can approximate the motion by many short segments,
write a difference equation, and then let the segment length tend to zero.
The result is an ordinary differential equation with a well-defined
continuum limit.

Quantum mechanics changes the sum over histories, not the conclusion.
Instead of selecting one classical path, we sum amplitudes over paths on the
surface.  Time may again be divided into small steps, and the limit is
equivalently described by a Schr\"odinger equation on the curved
configuration space.  For a smooth surface, the limit settles to a definite
propagator and spectrum.

\begin{figure}[H]
  \centering
  \includegraphics[width=0.72\textwidth]{lecture_09_figure_01.png}
  \caption{A smooth path on the surface and a jagged path from the
  time-sliced quantum sum.  Both remain intrinsic to the target surface.}
  \lectureframe{00:05:10}
  \label{fig:l09-particle-paths}
\end{figure}

\begin{classroomqa}
\audiencequestion{In the path integral, are all of the paths required to
remain inside the surface? \lecturetimestamp{00:05:39}}
\lecturerresponse{Yes.  The surface is the particle's configuration space,
not a sheet embedded in a larger room through which the particle may take a
shortcut.  Every intermediate point belongs to the surface, and all lengths
and actions are computed with its intrinsic metric.}
\end{classroomqa}

The corresponding question for a string is harder.  A string is not
specified only by a center-of-mass position.  It carries an infinite tower
of oscillators, so improving the resolution means admitting more dynamical
degrees of freedom.  The theory is consistent only if observables approach
a limit as that tower is restored.

\section{The Point Particle on a Sphere}

The sphere is the simplest curved test geometry.  Let its radius be \(R\),
and first place a particle of mass \(m\) on it.  For motion along a great
circle,
\begin{equation}
  E=\frac{1}{2}mv^2=\frac{p^2}{2m},
  \qquad
  L=pR.
  \label{eq:l09-particle-energy}
\end{equation}
Consequently,
\begin{equation}
  E=\frac{L^2}{2mR^2}
   =\frac{L^2}{2I},
  \qquad
  I=mR^2.
  \label{eq:l09-rotor-energy}
\end{equation}
The radius enters through the moment of inertia.  Quantizing the angular
momentum produces the familiar rigid-rotor spectrum, and nothing depends on
an ultraviolet cutoff.

\begin{figure}[H]
  \centering
  \includegraphics[width=0.74\textwidth]{lecture_09_figure_02.png}
  \caption{The point-particle control calculation:
  \(p^2/(2m)=L^2/(2mR^2)=L^2/(2I)\).}
  \lectureframe{00:11:25}
  \label{fig:l09-sphere-rotor}
\end{figure}

This untroubled example is important.  Curvature by itself is not the
problem.  The problem appears when curvature is coupled to the internal
quantum extension of a string.

\section{The Vacuum String Is Not Pointlike}

Consider one transverse coordinate of an open string and subtract its
center of mass.  The oscillator expansion used here is schematic:
\begin{equation}
  X(\sigma)-X_{\mathrm{cm}}
  \sim
  \sum_{n=1}^{\infty}
  \frac{a_n+a_n^\dagger}{\sqrt n}\cos(n\sigma).
  \label{eq:l09-mode-expansion}
\end{equation}
The essential ingredients are the cosine modes, the oscillator operators,
and the factor \(1/\sqrt n\).  Exact normalizations play no role in the
size argument.
\editorialnote{The precise expansion contains convention-dependent factors
of \(\alpha'\), numerical constants, and a choice of oscillator notation.
Equation \eqref{eq:l09-mode-expansion} retains only the dependence needed
for the ultraviolet estimate, just as in the board derivation.}

\begin{figure}[H]
  \centering
  \includegraphics[width=0.78\textwidth]{lecture_09_figure_03.png}
  \caption{The structural oscillator expansion and the double sum obtained
  on squaring it.}
  \lectureframe{00:15:20}
  \label{fig:l09-string-modes}
\end{figure}

The squared distance from the center of mass in the oscillator vacuum is
\begin{align}
  \expect{\bigl(X(\sigma)-X_{\mathrm{cm}}\bigr)^2}_0
  &\sim
  \sum_{n,m\geq1}
  \frac{\cos(n\sigma)\cos(m\sigma)}{\sqrt{nm}}\notag\\
  &\qquad\times
  \bra{0}(a_n+a_n^\dagger)\notag\\
  &\hspace{6em}\times(a_m+a_m^\dagger)\ket{0}.
  \label{eq:l09-size-double-sum}
\end{align}
Most terms vanish.  Two annihilation operators kill the ket; two creation
operators leave a state orthogonal to the bra; and a creation operator
followed by an annihilation operator also has zero vacuum expectation in
this ordering.  The surviving contraction is
\begin{equation}
  \bra{0}a_n a_m^\dagger\ket{0}=\delta_{nm}.
  \label{eq:l09-vacuum-contraction}
\end{equation}

\begin{figure}[H]
  \centering
  \includegraphics[width=0.74\textwidth]{lecture_09_figure_04.png}
  \caption{The surviving vacuum contraction collapses the double oscillator
  sum to a single sum.}
  \lectureframe{00:20:10}
  \label{fig:l09-vacuum-contraction}
\end{figure}

Thus
\begin{equation}
  \expect{\bigl(X(\sigma)-X_{\mathrm{cm}}\bigr)^2}_0
  \sim
  \sum_{n=1}^{\infty}\frac{\cos^2(n\sigma)}{n}.
  \label{eq:l09-size-single-sum}
\end{equation}
At a generic point along the string, or after averaging over \(\sigma\),
\(\cos^2(n\sigma)\) contributes roughly \(1/2\).  With a mode cutoff
\(N\),
\begin{equation}
  \expect{\bigl(X-X_{\mathrm{cm}}\bigr)^2}_{0,N}
  \sim
  \frac{1}{2}\sum_{n=1}^{N}\frac{1}{n}
  =\frac{1}{2}\log N+O(1).
  \label{eq:l09-logarithmic-size}
\end{equation}

\begin{figure}[H]
  \centering
  \includegraphics[width=0.72\textwidth]{lecture_09_figure_05.png}
  \caption{The board arrives at the harmonic sum controlling the
  ground-state spread.}
  \lectureframe{00:22:05}
  \label{fig:l09-harmonic-size}
\end{figure}

Keeping only, say, the first twenty-five modes gives a finite fuzzy object,
but it is not yet the continuum string.  Removing the cutoff causes the
mean-square spread to grow logarithmically.  This conclusion should not be
misread as a proof that flat-space string theory is inconsistent.  In flat
space the center-of-mass and oscillator motions separate, and delicate
cancellations in physical amplitudes preserve the expected locality.  The
growing size is a warning.  Curvature tells us whether the warning becomes
a failure.

\section{When Internal Spreading Changes the Geometry}

Place the string near the north pole of the sphere.  With every oscillator
turned off, it behaves like the point particle just studied.  As higher
modes are admitted, the string becomes a fluctuating tangle spread over the
cap.  Its center may still travel around the sphere, but the object whose
moment of inertia enters the rotational energy is now extended.

For an axis through the center of the sphere and the poles,
\begin{equation}
  I=\int r_\perp^2\,dm,
  \label{eq:l09-extended-inertia}
\end{equation}
where \(r_\perp\) is the perpendicular distance from a string element to
that axis.  A point at the equator of the orbit has the largest relevant
lever arm.  Portions of a string that droop toward the axis contribute less.
At fixed total mass this lowers the effective moment of inertia:
\begin{equation}
  I_{\mathrm{eff}}<mR^2.
  \label{eq:l09-inertia-reduction}
\end{equation}
The rotational energy therefore looks like that of a point object on a
smaller sphere,
\begin{equation}
  E_{\mathrm{rot}}
  =\frac{L^2}{2I_{\mathrm{eff}}}
  \equiv \frac{L^2}{2mR_{\mathrm{eff}}^2},
  \qquad R_{\mathrm{eff}}<R.
  \label{eq:l09-effective-radius}
\end{equation}

\begin{classroomqa}
\audiencequestion{Why should the moment of inertia tend toward zero if the
string is spreading over a larger region? \lecturetimestamp{00:33:23}}
\lecturerresponse{The statement is about the effective description under
successive coarse-graining, not a claim that the raw extended distribution
immediately has zero inertia.  The first band of modes turns the point into
a fuzzy lump that moves as though it lived on a smaller sphere.  The next
band resolves structure inside each such lump and repeats the same
renormalization.  Iterating the comparison drives the effective radius
inward.}
\end{classroomqa}

\begin{figure}[H]
  \centering
  \includegraphics[width=0.72\textwidth]{lecture_09_figure_06.png}
  \medskip

  \begin{adjustbox}{max width=\linewidth}
  \begin{tikzpicture}[x=1cm,y=1cm,>=Latex]
    \draw[thick] (0,0) circle (2.15);
    \draw[dashed] (0,-2.55) -- (0,2.55);
    \draw[blue!65!black,thick]
      (-0.22,2.05)
      .. controls (-1.22,1.75) and (-1.52,1.08) .. (-1.06,0.44)
      .. controls (-0.55,-0.10) and (0.35,0.02) .. (1.10,0.68)
      .. controls (1.48,1.12) and (1.16,1.80) .. (0.20,2.05);
    \draw[red!70!black,dashed] (0,0) circle (1.25);
    \draw[->] (2.70,1.55) -- (4.05,1.55);
    \node[align=center] at (3.37,1.98) {add shorter\\modes};
    \draw[thick] (6.15,0) circle (1.25);
    \draw[dashed] (6.15,-1.65) -- (6.15,1.65);
    \draw[blue!65!black,thick]
      (6.00,1.18)
      .. controls (5.42,0.88) and (5.46,0.25) .. (6.02,-0.06)
      .. controls (6.46,-0.28) and (6.90,0.45) .. (6.28,1.18);
    \node at (0,-3.00) {sphere of radius \(R\)};
    \node at (6.15,-2.05) {effective radius \(R_{\mathrm{eff}}\)};
  \end{tikzpicture}
  \end{adjustbox}
  \caption{The north-pole string on the board and a clean reconstruction of
  the effective-radius argument.  The dashed inner circle is an effective
  geometry, not a second physical shell inside the sphere.}
  \lectureframe{00:35:15}
  \label{fig:l09-shrinking-sphere}
\end{figure}

This is a renormalization-group statement.  A finite cutoff makes the string
a lump with finite internal resolution.  Raising the cutoff changes the
mechanics of that lump.  On a positively curved sphere each added shell of
modes makes the background inferred from the long-distance motion smaller.
If the result never stabilizes, there is no cutoff-independent string theory
on that background.

\section{The Metric Must Reach a Fixed Point}

Now replace the sphere by a general target-space metric
\(g_{\mu\nu}(x)\).  Integrating in a thin logarithmic shell of oscillator
modes changes that metric by a local symmetric rank-two tensor.  The change
must vanish in flat space and begin with two derivatives of the metric.
At leading order, covariance leaves the Ricci tensor as the relevant
candidate:
\begin{equation}
  \frac{\partial g_{\mu\nu}}{\partial\log\Lambda}
  =-C\,R_{\mu\nu}+\cdots ,
  \qquad C>0
  \label{eq:l09-ricci-flow}
\end{equation}
in the direction of coarse-graining used in the geometric argument.
Positive curvature then drives a sphere inward, as required.
\editorialnote{The normalization and sign depend on whether the
renormalization parameter increases toward the ultraviolet or toward the
infrared.  In conventional worldsheet notation the leading metric beta
function is \(\beta^g_{\mu\nu}=\alpha'R_{\mu\nu}+\cdots\).  The invariant
statement used here is that cutoff independence requires the beta function,
and hence \(R_{\mu\nu}\) at this order, to vanish.}

\begin{figure}[H]
  \centering
  \includegraphics[width=0.76\textwidth]{lecture_09_figure_07.png}
  \caption{The metric change is written schematically as minus the Ricci
  tensor, with a smoothing curve drawn beside it.}
  \lectureframe{00:43:00}
  \label{fig:l09-ricci-flow-board}
\end{figure}

\begin{classroomqa}
\audiencequestion{Where does the oscillator cutoff \(N\) appear in this
flow equation? \lecturetimestamp{00:43:38}}
\lecturerresponse{It is the scale variable.  The string-size sum changes
logarithmically when the highest included mode changes, so the natural
evolution parameter is \(\log N\), or equivalently \(\log\Lambda\).
Equation \eqref{eq:l09-ricci-flow} describes the infinitesimal change caused
by one logarithmic shell of modes.}
\end{classroomqa}

An admissible background is a fixed point of this flow.  To leading order,
\begin{equation}
  \boxed{R_{\mu\nu}=0.}
  \label{eq:l09-ricci-flat}
\end{equation}
Ricci flat does not mean that every component of the Riemann tensor
vanishes.  Gravitational waves, for example, can have nonzero curvature in
vacuum while satisfying \(R_{\mu\nu}=0\).

\subsection{The Same Condition Is the Vacuum Einstein Equation}

With conventional constants suppressed, Einstein's equation is
\begin{equation}
  G_{\mu\nu}
  \equiv
  R_{\mu\nu}-\frac12 g_{\mu\nu}R
  =T_{\mu\nu}.
  \label{eq:l09-einstein}
\end{equation}

\begin{figure}[H]
  \centering
  \includegraphics[width=0.78\textwidth]{lecture_09_figure_08.png}
  \caption{The Einstein tensor and stress tensor on the board.}
  \lectureframe{00:47:15}
  \label{fig:l09-einstein-board}
\end{figure}

In vacuum \(T_{\mu\nu}=0\).  Taking the trace in \(D\) spacetime
dimensions gives
\begin{equation}
  0=g^{\mu\nu}G_{\mu\nu}
   =\left(1-\frac{D}{2}\right)R.
  \label{eq:l09-einstein-trace}
\end{equation}
For \(D\neq2\), \(R=0\), and substitution back into
\eqref{eq:l09-einstein} yields \(R_{\mu\nu}=0\).  Thus the condition that
the string lose its cutoff dependence is the vacuum gravitational field
equation.  Gravity has not been appended to the theory; it has appeared as
the price of consistency.

The word \emph{vacuum} does not mean that nothing happens.  It means there
is no stress-energy source in the region.  Source-free Maxwell equations
admit electromagnetic waves, and source-free Einstein equations admit
gravitational waves.  The geometry may be dynamical even when its Ricci
tensor vanishes.

\begin{classroomqa}
\audiencequestion{Is this condition connected with conformal invariance?
\lecturetimestamp{00:52:20}}
\lecturerresponse{Yes, although the question is partly lost in the
recording.  The relevant symmetry is conformal invariance of the
\emph{worldsheet} quantum theory.  Requiring the worldsheet scale anomaly
to vanish is what sets the metric beta function to zero.  It should not be
confused with demanding conformal invariance of spacetime itself.}
\end{classroomqa}

There is a second consistency condition.  Worldsheet conformal invariance
also fixes the critical spacetime dimension: ten for superstrings and
twenty-six for the bosonic string.  Ricci flatness alone is therefore not
the complete list of requirements, but it is the part that displays the
Einstein equation most directly.

\section{Compact Dimensions Rather Than Missing Dimensions}

A superstring background has nine spatial dimensions and one time
dimension.  Our macroscopic world displays only three large spatial
directions, so six spatial dimensions must be hidden or otherwise rendered
inaccessible at ordinary scales.  Compactification keeps them as genuine
directions while giving them small finite extent.

\begin{classroomqa}
\audiencequestion{If different observers carry different clocks, does that
amount to having more than one time dimension? \lecturetimestamp{00:56:59}}
\lecturerresponse{No.  Many clocks are many instruments measuring
coordinates in one spacetime; they do not create independent timelike
directions.  The extra dimensions considered here are spatial.  The
familiar quantum and causal constructions do not supply a comparable
consistent theory with several local time directions, so that possibility
is not part of this compactification model.}
\end{classroomqa}

Imagine creatures who believe their world is a line.  At low resolution
they can move only forward and backward.  Under a sufficiently powerful
microscope, however, the line is revealed to be a very thin cylinder.  A
small creature can move around its circular direction as well as along its
length.  The second direction was always present; its circumference was too
small for coarse probes to resolve.

\begin{figure}[H]
  \centering
  \includegraphics[width=0.78\textwidth]{lecture_09_figure_09.png}
  \caption{A narrow cylinder and trajectories that wind around its hidden
  circular direction.}
  \lectureframe{01:01:00}
  \label{fig:l09-hidden-cylinder}
\end{figure}

A cylinder can be represented as a strip whose upper and lower edges are
identified.  Crossing the upper edge returns through the lower edge at the
same longitudinal position.  There is no physical boundary at either edge.
The same operation can be performed with two directions: identify the left
and right edges of a rectangle and also its top and bottom edges.  The
result is a two-torus.

\begin{figure}[H]
  \centering
  \includegraphics[width=0.76\textwidth]{lecture_09_figure_10.png}
  \medskip

  \begin{tikzpicture}[x=0.82cm,y=0.82cm,>=Latex]
    \draw[thick] (0,0) rectangle (4,2.15);
    \draw[blue!65!black,->] (0,2.48) -- (4,2.48);
    \draw[blue!65!black,->] (4,-0.34) -- (0,-0.34);
    \node at (2,-0.82) {one periodic direction};
    \draw[thick] (6.1,0) rectangle (10.1,2.15);
    \draw[blue!65!black,->] (6.1,2.48) -- (10.1,2.48);
    \draw[blue!65!black,->] (10.1,-0.34) -- (6.1,-0.34);
    \draw[red!70!black,->] (5.76,0) -- (5.76,2.15);
    \draw[red!70!black,->] (10.44,2.15) -- (10.44,0);
    \node at (8.1,-0.82) {two periodic directions};
  \end{tikzpicture}
  \caption{The board's slab and a clean fundamental-domain description.
  Identifying one edge pair gives a cylinder; identifying both pairs gives
  a torus.}
  \lectureframe{01:04:40}
  \label{fig:l09-periodic-identification}
\end{figure}

The torus is a topology before it is a bagel-shaped surface in
three-dimensional space.  Cut a torus once and it becomes a cylinder; cut
again and it becomes a rectangle, provided that the edge identifications
are remembered.  A flat metric may be put on this rectangle, so the
resulting flat torus has zero Riemann curvature locally and is certainly
Ricci flat.

A straight trajectory on the rectangle continues straight across an
identified edge.  If its slope is a rational ratio of the two periods, the
trajectory eventually closes.  If the ratio is irrational, it never closes
and becomes dense on the torus.

\begin{figure}[H]
  \centering
  \includegraphics[width=0.76\textwidth]{lecture_09_figure_11.png}
  \caption{Repeated straight segments in a fundamental domain represent
  one uninterrupted trajectory on the torus.}
  \lectureframe{01:11:00}
  \label{fig:l09-torus-trajectory}
\end{figure}

Higher-dimensional tori are formed in the same way.  A cube with each pair
of opposite faces identified is a three-torus, and six compact directions
can be modeled by a six-torus.  One should not try to embed all of these
objects in ordinary three-dimensional space; the face-identification
description is intrinsic and complete.

\begin{classroomqa}
\audiencequestion{Must an extra dimension necessarily come back on itself?
\lecturetimestamp{01:15:03}}
\lecturerresponse{Not as a theorem.  It could be noncompact, but then it
would normally appear as another large direction in the observed
long-distance world.  Compactness is the simplest way to reconcile extra
dimensions with their apparent absence.  Why nature chooses a particular
compactification is a separate dynamical question.}
\end{classroomqa}

\begin{classroomqa}
\audiencequestion{Why assume the unseen dimensions are ordinary spatial
directions curled up in this way rather than something altogether
different? \lecturetimestamp{01:15:51}}
\lecturerresponse{Because this is a definite solvable hypothesis, not
because it has been proved to be the unique description of nature.
Compact spatial manifolds let us construct consistent examples and compute
their spectra.  More exotic possibilities must earn their place by doing
the same.}
\end{classroomqa}

\begin{classroomqa}
\audiencequestion{Distances on the familiar bagel-shaped torus do not look
like distances in a rectangle.  Which geometry is meant?
\lecturetimestamp{01:18:11}}
\lecturerresponse{The term \emph{torus} fixes the identifications, not one
unique metric.  The embedded bagel inherits a curved metric from
three-dimensional space.  The periodically identified rectangle carries a
flat metric.  They have the same topology and different geometry; for the
simple compactification we mean the flat one.}
\end{classroomqa}

\section{The Shape Parameters of a Flat Torus}

Even a flat two-torus is not unique.  Its fundamental parallelogram has an
overall area, an aspect ratio, and a shear.  Changing these parameters
changes distances and therefore changes the spectrum, without changing the
topology.

\begin{figure}[H]
  \centering
  \includegraphics[width=0.74\textwidth]{lecture_09_figure_12.png}
  \medskip

  \begin{tikzpicture}[x=0.9cm,y=0.9cm,>=Latex]
    \draw[thick] (0,0) rectangle (3.4,2.1);
    \draw[<->] (0,-0.35) -- node[below] {\(L_1\)} (3.4,-0.35);
    \draw[<->] (-0.35,0) -- node[left] {\(L_2\)} (-0.35,2.1);
    \draw[thick] (5.4,0) -- (8.8,0) -- (9.7,2.1) -- (6.3,2.1) -- cycle;
    \draw[dashed] (5.4,0) -- (6.3,2.1);
    \draw[->] (5.75,2.48) -- node[above] {shear} (6.45,2.48);
    \node at (1.7,-1.00) {area and aspect ratio};
    \node at (7.55,-1.00) {complex structure};
  \end{tikzpicture}
  \caption{The rectangular and sheared fundamental domains used to discuss
  torus moduli.}
  \lectureframe{01:23:00}
  \label{fig:l09-torus-moduli}
\end{figure}

For a two-torus it is useful to package the shape into a complex number
\(\tau=\tau_1+i\tau_2\): the real part records shear and the imaginary part
records an aspect ratio after an overall scale is separated.  Interchanging
the two basic cycles or changing their integer basis can make apparently
different parallelograms describe the same torus.

\begin{classroomqa}
\audiencequestion{What happens if one edge is flipped before the
identification? \lecturetimestamp{01:24:20}}
\lecturerresponse{That produces a different topology.  Identifying with a
reflection gives a nonorientable surface, the Klein bottle, rather than an
ordinary torus.}
\end{classroomqa}

\begin{classroomqa}
\audiencequestion{What if the twist is not uniform along the edge?
\lecturetimestamp{01:25:01}}
\lecturerresponse{A smooth nonuniform deformation can be transferred into
the metric on the fundamental domain.  The topological identification and
the geometrical data must be kept distinct: changing the gluing may change
topology, while smoothly changing distances changes the metric.}
\end{classroomqa}

\editorialnote{For a flat two-torus, the area is the simplest Kahler
modulus and \(\tau\) is the complex-structure modulus.  A continuous shear
is part of \(\tau\).  A Dehn twist is instead a discrete mapping-class
operation, represented by \(\tau\mapsto\tau+1\); it is not itself the
continuous modulus.  This sharpens the intentionally informal terminology
used during the board discussion.}

Flat tori are useful because they solve the consistency condition
immediately, but their symmetry is too generous for a realistic model of
particle physics.  More intricate Ricci-flat spaces, notably Calabi--Yau
manifolds, provide richer compactifications.  Their detailed construction
lies beyond this discussion; the principle remains the same.

\section{Reading a Compact Circle from Its Spectrum}

After the break, reduce the compact space to one circle of circumference
\(2\pi R\).  A wavefunction must be single-valued after one circuit:
\begin{equation}
  e^{ip_c(x+2\pi R)/\hbar}=e^{ip_cx/\hbar}.
\end{equation}
Hence
\begin{equation}
  p_cR=n\hbar,
  \qquad
  p_c=\frac{n\hbar}{R},
  \qquad n\in\mathbb Z.
  \label{eq:l09-kk-momentum}
\end{equation}

\begin{figure}[H]
  \centering
  \includegraphics[width=0.74\textwidth]{lecture_09_figure_13.png}
  \caption{The compact circumference \(2\pi R\) and quantized momentum
  \(p_c=n\hbar/R\).}
  \lectureframe{01:32:15}
  \label{fig:l09-kk-board}
\end{figure}

For a massless particle in the full spacetime,
\begin{equation}
  E=c|\mathbf P|.
  \label{eq:l09-massless-dispersion}
\end{equation}

\begin{figure}[H]
  \centering
  \includegraphics[width=0.70\textwidth]{lecture_09_figure_14.png}
  \caption{The massless dispersion relation used to reinterpret compact
  momentum as lower-dimensional mass.}
  \lectureframe{01:33:45}
  \label{fig:l09-massless-board}
\end{figure}

Split the momentum into noncompact and compact parts and set
\(c=\hbar=1\):
\begin{equation}
  E^2=\mathbf p_{\mathrm{large}}^2+\frac{n^2}{R^2}.
  \label{eq:l09-kk-dispersion}
\end{equation}
An observer who cannot resolve the compact direction compares this with
\(E^2=\mathbf p_{\mathrm{large}}^2+m^2\) and infers a tower of masses
\begin{equation}
  m_n=\frac{|n|}{R}.
  \label{eq:l09-kk-masses}
\end{equation}
The two signs of \(n\) describe opposite compact momenta with the same
positive energy, not positive and negative masses.  The spacing \(1/R\)
therefore reveals the radius spectroscopically.  This is the
Kaluza--Klein mechanism: momentum in a hidden direction appears as mass in
the visible dimensions.

\section{Winding States Complete the Spectrum}

A closed string has a possibility unavailable to a point particle: it can
wrap around the circle.  If the winding number is \(w\), the string closes
only after covering the compact direction \(w\) times.  Its length is
\(2\pi |w|R\), so tension supplies an energy
\begin{equation}
  m_w
  =2\pi T\,|w|R
  =\frac{|w|R}{\alpha'}.
  \label{eq:l09-winding-mass}
\end{equation}
The lecture chooses units with \(\alpha'=1\), reducing this to
\begin{equation}
  m_w=|w|R.
  \label{eq:l09-winding-lecture-units}
\end{equation}
Positive and negative \(w\) are opposite orientations around the circle.

\begin{figure}[H]
  \centering
  \includegraphics[width=0.76\textwidth]{lecture_09_figure_15.png}
  \caption{A closed string wound around the compact cylinder, with the
  compact radius marked on the board.}
  \lectureframe{01:39:45}
  \label{fig:l09-winding-board}
\end{figure}

The compact circle thus produces two complementary ladders:
\begin{equation}
  m_n=\frac{|n|}{R},
  \qquad
  m_w=\frac{|w|R}{\alpha'}.
  \label{eq:l09-two-towers}
\end{equation}
When \(R\) is large, momentum levels are closely spaced and winding states
are heavy.  When \(R\) is small, momentum levels are widely spaced and
winding states are light.

\begin{figure}[H]
  \centering
  \begin{tikzpicture}[x=1cm,y=0.9cm,>=Latex]
    \draw[->] (0,0) -- (0,4.7) node[above] {\(m\)};
    \foreach \y/\lab in {0/0,1.1/{1/R},2.2/{2/R},3.3/{3/R},4.4/{4/R}} {
      \draw (-0.16,\y) -- (0.16,\y);
      \node[left] at (-0.22,\y) {\(\lab\)};
    }
    \node at (0,-0.62) {momentum};
    \draw[<->,very thick,blue!65!black] (2.05,2.15) -- (4.05,2.15)
      node[midway,above] {\(R\leftrightarrow\alpha'/R\)};
    \draw[->] (6.1,0) -- (6.1,4.7) node[above] {\(m\)};
    \foreach \y/\lab in {0/0,1.1/{R/\alpha'},2.2/{2R/\alpha'},3.3/{3R/\alpha'},4.4/{4R/\alpha'}} {
      \draw (5.94,\y) -- (6.26,\y);
      \node[right] at (6.32,\y) {\(\lab\)};
    }
    \node at (6.1,-0.62) {winding};
  \end{tikzpicture}
  \caption{Momentum and winding towers exchange when the radius is
  inverted in string units.}
  \label{fig:l09-dual-mass-towers}
\end{figure}

\section{T-Duality: Small and Large Circles Are the Same Physics}

The two towers are interchanged by
\begin{equation}
  R\longleftrightarrow\frac{\alpha'}{R},
  \qquad
  n\longleftrightarrow w.
  \label{eq:l09-t-duality}
\end{equation}
In the units used on the board, \(\alpha'=1\), so this is simply
\(R\leftrightarrow1/R\).  At the self-dual radius
\(R=\sqrt{\alpha'}\), momentum and winding scales coincide.

\begin{figure}[H]
  \centering
  \includegraphics[width=0.74\textwidth]{lecture_09_figure_16.png}
  \caption{The sparse and dense mass ladders compared on the board as the
  radius crosses the self-dual scale.}
  \lectureframe{01:43:00}
  \label{fig:l09-duality-board}
\end{figure}

The exchange of the elementary mass formulas is only the visible edge of
the result.  Closed-string oscillator states and scattering amplitudes are
also mapped consistently.  A circle smaller than the string scale is not a
new increasingly tiny geometry; it has an equivalent description as a
larger dual circle.  Point particles cannot make this identification
because they possess momentum modes but no winding modes.

\begin{classroomqa}
\audiencequestion{Is this merely the same as exchanging a vertical and a
horizontal direction? \lecturetimestamp{01:44:38}}
\lecturerresponse{No.  It exchanges two different quantum numbers:
momentum around the circle and the number of times a closed string wraps
the circle.  The radius must be inverted at the same time.}
\end{classroomqa}

\begin{classroomqa}
\audiencequestion{Do the descriptions at \(R<\sqrt{\alpha'}\) and
\(R>\sqrt{\alpha'}\) refer only to bosons, or also to fermions?
\lecturetimestamp{01:46:30}}
\lecturerresponse{The duality is a symmetry of the full compactified string
theory, not just the simple bosonic mass ladder used to expose it.  The
detailed mapping of sectors depends on the theory, but the small--large
equivalence is not restricted to one particle species.}
\end{classroomqa}

\begin{classroomqa}
\audiencequestion{Are the two states being compared open strings or closed
strings? \lecturetimestamp{01:47:32}}
\lecturerresponse{They are closed-string states.  A freely propagating
closed string can wind a compact circle.  An ordinary open string with free
endpoints does not carry the same conserved winding number.}
\end{classroomqa}
\editorialnote{In the fuller open-string account, T-duality exchanges
Neumann and Dirichlet boundary conditions and introduces D-branes.  That is
why open-string T-duality is richer than simply copying the closed-string
winding formula.}

\begin{classroomqa}
\audiencequestion{Is the interchange analogous to the exchange of space and
time near a black-hole horizon? \lecturetimestamp{01:48:00}}
\lecturerresponse{No.  Horizon coordinates and causal character are a
different issue.  Here the two descriptions concern one compact spatial
circle and the exchange of its momentum and winding sectors.}
\end{classroomqa}

\begin{classroomqa}
\audiencequestion{Does this cure the short-distance self-energy problem by
giving a minimum length? \lecturetimestamp{01:48:35}}
\lecturerresponse{It shows that radii below the string scale do not define
new inequivalent circle backgrounds: decreasing \(R\) beyond the self-dual
point increases the dual radius.  This is a precise stringy obstruction to
probing arbitrarily small circle geometry, although it should not be turned
into a claim that every conceivable notion of distance has one universal
hard lower bound.}
\end{classroomqa}

\section{Why Winding Sectors Cannot Be Omitted}

\begin{classroomqa}
\audiencequestion{Could we define the theory using only unwound strings and
leave winding states out? \lecturetimestamp{01:48:56}}
\lecturerresponse{No.  Closed-string interactions generate them.  Stretch
an unwound loop around the compact direction so that one segment follows
the circle one way and another follows it the opposite way.  The net
winding is still zero.  A reconnection can then split the loop into two
closed strings with winding \(+1\) and \(-1\).  The pair may separate, so
winding states must belong to the spectrum even if the initial state had
none.}
\end{classroomqa}

\begin{figure}[H]
  \centering
  \includegraphics[width=0.74\textwidth]{lecture_09_figure_17.png}
  \caption{The stretched zero-winding loop arranged around the compact
  direction before reconnection.}
  \lectureframe{01:50:30}
  \label{fig:l09-winding-pair-before}
\end{figure}

\begin{figure}[H]
  \centering
  \includegraphics[width=0.72\textwidth]{lecture_09_figure_18.png}
  \medskip

  \begin{tikzpicture}[x=0.9cm,y=0.9cm,>=Latex]
    \draw[thick] (0,0) ellipse (1.0 and 1.55);
    \draw[blue!65!black,thick]
      (0.08,1.35) .. controls (2.0,1.15) and (2.0,-1.15) .. (0.08,-1.35);
    \draw[red!70!black,thick]
      (-0.08,-1.35) .. controls (-2.0,-1.15) and (-2.0,1.15) .. (-0.08,1.35);
    \node at (0,-2.15) {net \(w=0\)};
    \draw[->,thick] (2.45,0) -- (3.75,0) node[midway,above] {split};
    \draw[thick] (5.05,0.62) ellipse (0.68 and 1.02);
    \draw[blue!65!black,very thick,->]
      (5.05,1.64) arc[start angle=90,end angle=-250,x radius=0.68,y radius=1.02];
    \node at (5.05,-0.82) {\(w=+1\)};
    \draw[thick] (7.45,0.62) ellipse (0.68 and 1.02);
    \draw[red!70!black,very thick,<-]
      (7.45,1.64) arc[start angle=90,end angle=-250,x radius=0.68,y radius=1.02];
    \node at (7.45,-0.82) {\(w=-1\)};
  \end{tikzpicture}
  \caption{The board's separated wound loop and a clean reconstruction of
  winding--antiwinding pair production.  Total winding remains zero.}
  \lectureframe{01:51:30}
  \label{fig:l09-winding-pair-after}
\end{figure}

\begin{classroomqa}
\audiencequestion{Does producing the winding pair require extra energy?
\lecturetimestamp{01:52:11}}
\lecturerresponse{Yes.  Stretching a string costs tension energy, and the
final wound strings have the masses in
\eqref{eq:l09-winding-mass}.  A sufficiently energetic collision can
produce the pair, just as a collision can produce a particle and its
antiparticle while preserving the relevant conserved charge.}
\end{classroomqa}

\section{Can One Hear the Shape of a String Background?}

For an ordinary compactified particle theory, the spacing of the
Kaluza--Klein spectrum lets us infer the radius: one tries to hear the shape
of the compact space from its normal modes.  String theory makes the inverse
problem subtler.  Momentum modes hear \(1/R\), winding modes hear
\(R/\alpha'\), and the complete spectrum is unchanged when the two are
exchanged.  Two radii that point-particle geometry regards as different are
one string background described in dual variables.

The lecture has therefore arrived at both of its promised conclusions by
one route.  Ultraviolet consistency forces the target metric to satisfy the
vacuum Einstein equation, so gravity is unavoidable.  Compactification then
turns hidden geometry into particle spectra, but winding states prevent that
spectrum from assigning a unique classical radius below the string scale.
Consistency does not merely select which geometries are allowed; through
T-duality, it changes what counts as a distinct geometry at all.
